% ===========================================================================
% Plantilla LaTeX para tablas de la tesis
% Uso: % ===========================================================================
% Plantilla LaTeX para tablas de la tesis
% Uso: % ===========================================================================
% Plantilla LaTeX para tablas de la tesis
% Uso: % ===========================================================================
% Plantilla LaTeX para tablas de la tesis
% Uso: \input{../tables/template_tabla.tex} (adaptar contenido)
% Requiere en el preámbulo: \usepackage{booktabs, siunitx, multirow, array}
% ===========================================================================

% Ejemplo de tabla con booktabs y alineación numérica
\begin{table}[htbp]
  \centering
  \caption{Descripción de la tabla en tiempo pasivo y descriptivo.}
  \label{tab:etiqueta}
  \begin{tabular}{@{}lS[table-format=3.2]S[table-format=2.2]c@{}}
    \toprule
    {Parámetro} & {Valor 1} & {Valor 2} & {Unidad} \\
    \midrule
    Precisión media & 101.5 & 14.78 & \si{\micro\second} \\
    Offset estabilizado & 1.265 & 2.555 & \si{\micro\second} \\
    Tiempo convergencia & 4.0 & 0.5 & \si{\second} \\
    \bottomrule
  \end{tabular}
\end{table}

% Notas:
% - 'S' alinea números por el punto decimal vía siunitx.
% - @{} elimina padding lateral para tablas más compactas.
% - \toprule, \midrule, \bottomrule en lugar de \hline (estilo booktabs).
% - Sin líneas verticales.
 (adaptar contenido)
% Requiere en el preámbulo: \usepackage{booktabs, siunitx, multirow, array}
% ===========================================================================

% Ejemplo de tabla con booktabs y alineación numérica
\begin{table}[htbp]
  \centering
  \caption{Descripción de la tabla en tiempo pasivo y descriptivo.}
  \label{tab:etiqueta}
  \begin{tabular}{@{}lS[table-format=3.2]S[table-format=2.2]c@{}}
    \toprule
    {Parámetro} & {Valor 1} & {Valor 2} & {Unidad} \\
    \midrule
    Precisión media & 101.5 & 14.78 & \si{\micro\second} \\
    Offset estabilizado & 1.265 & 2.555 & \si{\micro\second} \\
    Tiempo convergencia & 4.0 & 0.5 & \si{\second} \\
    \bottomrule
  \end{tabular}
\end{table}

% Notas:
% - 'S' alinea números por el punto decimal vía siunitx.
% - @{} elimina padding lateral para tablas más compactas.
% - \toprule, \midrule, \bottomrule en lugar de \hline (estilo booktabs).
% - Sin líneas verticales.
 (adaptar contenido)
% Requiere en el preámbulo: \usepackage{booktabs, siunitx, multirow, array}
% ===========================================================================

% Ejemplo de tabla con booktabs y alineación numérica
\begin{table}[htbp]
  \centering
  \caption{Descripción de la tabla en tiempo pasivo y descriptivo.}
  \label{tab:etiqueta}
  \begin{tabular}{@{}lS[table-format=3.2]S[table-format=2.2]c@{}}
    \toprule
    {Parámetro} & {Valor 1} & {Valor 2} & {Unidad} \\
    \midrule
    Precisión media & 101.5 & 14.78 & \si{\micro\second} \\
    Offset estabilizado & 1.265 & 2.555 & \si{\micro\second} \\
    Tiempo convergencia & 4.0 & 0.5 & \si{\second} \\
    \bottomrule
  \end{tabular}
\end{table}

% Notas:
% - 'S' alinea números por el punto decimal vía siunitx.
% - @{} elimina padding lateral para tablas más compactas.
% - \toprule, \midrule, \bottomrule en lugar de \hline (estilo booktabs).
% - Sin líneas verticales.
 (adaptar contenido)
% Requiere en el preámbulo: \usepackage{booktabs, siunitx, multirow, array}
% ===========================================================================

% Ejemplo de tabla con booktabs y alineación numérica
\begin{table}[htbp]
  \centering
  \caption{Descripción de la tabla en tiempo pasivo y descriptivo.}
  \label{tab:etiqueta}
  \begin{tabular}{@{}lS[table-format=3.2]S[table-format=2.2]c@{}}
    \toprule
    {Parámetro} & {Valor 1} & {Valor 2} & {Unidad} \\
    \midrule
    Precisión media & 101.5 & 14.78 & \si{\micro\second} \\
    Offset estabilizado & 1.265 & 2.555 & \si{\micro\second} \\
    Tiempo convergencia & 4.0 & 0.5 & \si{\second} \\
    \bottomrule
  \end{tabular}
\end{table}

% Notas:
% - 'S' alinea números por el punto decimal vía siunitx.
% - @{} elimina padding lateral para tablas más compactas.
% - \toprule, \midrule, \bottomrule en lugar de \hline (estilo booktabs).
% - Sin líneas verticales.
